\documentclass[aspectratio=169,12pt]{beamer}
% ============================================================================
% THE SOVEREIGN LEDGER — reworked edition, 5 August 2026
% Compile with: latexmk -xelatex sovereign-ledger.tex
%
% TIMING PLAN (~90 minutes total; per-section timings in section-frame notes):
%   Framing 12' | I 10' | II 10' | III 11' | IV 12' | V 12' | VI 11'
%   VII 9' | VIII 12' | Close 1'   — appendix not presented.
% Slides marked "% FAST" are 30–45 second glance slides.
%
% SPEAKER NOTES: every frame carries a \note. To render them, uncomment ONE of:
% \setbeameroption{show notes on second screen=right}   % dual-screen (pgfpages)
% \setbeameroption{show notes}                          % interleaved pages
% ============================================================================

\usepackage{fontspec}
\setsansfont{Latin Modern Sans}
\usepackage{microtype}
\usepackage{booktabs}
\usepackage{tabularx}
\usepackage{array}
\usepackage{ragged2e}
\usepackage{graphicx}
\usepackage{hyperref}
\usepackage{pgfpages} % required only if a show-notes option above is enabled

\definecolor{Ink}{HTML}{17242D}
\definecolor{Paper}{HTML}{F5F1E8}
\definecolor{Rust}{HTML}{B94E35}
\definecolor{Ocean}{HTML}{2F6B84}
\definecolor{Sage}{HTML}{6D806F}
\definecolor{Gold}{HTML}{B68A3E}
\definecolor{Mist}{HTML}{DDE5E6}
\definecolor{Soft}{HTML}{E9E3D8}
\definecolor{Quiet}{HTML}{5E6B73}
\definecolor{White}{HTML}{FFFFFF}

\hypersetup{colorlinks=true,urlcolor=Ocean,linkcolor=Ink,citecolor=Ocean}

\setbeamercolor{normal text}{fg=Ink,bg=Paper}
\setbeamercolor{structure}{fg=Rust}
\setbeamercolor{frametitle}{fg=Ink,bg=Paper}
\setbeamercolor{title}{fg=White}
\setbeamercolor{subtitle}{fg=Mist}
\setbeamercolor{footline}{fg=Quiet,bg=Paper}
\setbeamercolor{block title}{fg=White,bg=Ocean}
\setbeamercolor{block body}{fg=Ink,bg=Mist}

\setbeamerfont{title}{size=\fontsize{30}{33},series=\bfseries}
\setbeamerfont{subtitle}{size=\fontsize{13.5}{16}}
\setbeamerfont{frametitle}{size=\fontsize{21}{24},series=\bfseries}
\setbeamerfont{normal text}{size=\fontsize{11.2}{13.5}}
\setbeamerfont{itemize/enumerate body}{size=\fontsize{11.2}{13.5}}
\setbeamerfont{itemize/enumerate subbody}{size=\fontsize{9.5}{11.5}}

\setbeamertemplate{navigation symbols}{}
\setbeamertemplate{itemize item}{\raisebox{0.15ex}{\color{Rust}\rule{0.55ex}{0.55ex}}}
\setbeamertemplate{itemize subitem}{\raisebox{0.15ex}{\color{Ocean}\rule{0.42ex}{0.42ex}}}
\setbeamersize{text margin left=0.62cm,text margin right=0.62cm}

\setbeamertemplate{frametitle}{%
  \vspace{0.22cm}%
  \begin{beamercolorbox}[wd=\paperwidth,leftskip=0.62cm,rightskip=0.62cm]{frametitle}%
    \usebeamerfont{frametitle}\hyphenpenalty=10000\exhyphenpenalty=10000 \insertframetitle\par
    \vspace{0.12cm}\color{Rust}\rule{1.25cm}{1.6pt}
  \end{beamercolorbox}%
}

\setbeamertemplate{footline}{%
  \leavevmode%
  \hbox{%
    \begin{beamercolorbox}[wd=.78\paperwidth,ht=2.65ex,dp=1.1ex,leftskip=.62cm]{footline}%
      \fontsize{6}{7}\selectfont THE SOVEREIGN LEDGER \quad | \quad 5 AUGUST 2026
    \end{beamercolorbox}%
    \begin{beamercolorbox}[wd=.22\paperwidth,ht=2.65ex,dp=1.1ex,rightskip=.62cm plus1fil]{footline}%
      \hfill\fontsize{6}{7}\selectfont \insertframenumber~/~\inserttotalframenumber
    \end{beamercolorbox}%
  }%
}

\newcommand{\rust}[1]{\textcolor{Rust}{\textbf{#1}}}
\newcommand{\blue}[1]{\textcolor{Ocean}{\textbf{#1}}}
\newcommand{\muted}[1]{\textcolor{Quiet}{#1}}
\newcommand{\claim}[1]{\begin{center}\fontsize{16}{19}\selectfont\bfseries #1\end{center}}
\newcommand{\source}[1]{\vfill\par\fontsize{5.4}{6.5}\selectfont\color{Quiet}Sources: #1}
\newcommand{\smallsource}[1]{\par\fontsize{5}{6}\selectfont\color{Quiet}Sources: #1}
\newcommand{\tightitems}{\setlength{\itemsep}{0.42em}\setlength{\parskip}{0pt}}
\newcommand{\verytightitems}{\setlength{\itemsep}{0.22em}\setlength{\parskip}{0pt}}
\newcolumntype{Y}{>{\RaggedRight\arraybackslash}X}
\newcolumntype{L}[1]{>{\RaggedRight\arraybackslash}p{#1}}

% Layer tag (refers to the eight-layer control framework)
\newcommand{\lyr}[1]{\textcolor{Gold}{\fontsize{7.3}{8.4}\selectfont\bfseries L#1}}
% Evidence tiers (defined on the roadmap slide; used in all assessments)
\newcommand{\tdoc}{\textcolor{Sage}{\bfseries documented}}
\newcommand{\tinf}{\textcolor{Ocean}{\bfseries inference}}
\newcommand{\tunp}{\textcolor{Quiet}{\bfseries unproven}}

\newenvironment{statement}{%
  \begin{beamercolorbox}[sep=.18cm,rounded=false,wd=\linewidth]{block body}%
  \fontsize{13}{15.5}\selectfont\bfseries
}{\end{beamercolorbox}}

% \sectionframe[strapline]{kicker}{title}{timing note}
\newcommand{\sectionframe}[4][Payments, clearing, title and programmable finance]{%
  \begingroup
  \setbeamercolor{background canvas}{bg=Ink}
  \begin{frame}[plain]
    \vspace{1.05cm}
    {\color{Gold}\fontsize{10}{12}\selectfont\bfseries #2\par}
    \vspace{0.36cm}
    {\color{White}\fontsize{26}{30}\selectfont\bfseries\hyphenpenalty=10000 #3\par}
    \vfill
    {\color{Mist}\rule{2.1cm}{2pt}}\par
    \vspace{0.28cm}
    {\color{Mist}\fontsize{8.5}{10.5}\selectfont #1}
    \note{[Timing] #4. [Sources] Section transition; no external claim.}
  \end{frame}
  \endgroup
}

\title{Project Agor\'a}
\subtitle{Payments, clearing and securities registries in the programmable transition}
\date{5 August 2026}

\begin{document}

% ============================ TITLE ============================
\begingroup
\setbeamercolor{background canvas}{bg=Ink}
\begin{frame}[plain]
  \vspace{0.62cm}
  {\color{Gold}\fontsize{10}{12}\selectfont\bfseries CONTROL, DEPENDENCE AND CAPTURE\par}
  \vspace{0.28cm}
  {\usebeamerfont{title}\color{White}\inserttitle\par}
  \vspace{0.32cm}
  {\usebeamerfont{subtitle}\color{Mist}\insertsubtitle\par}
  \vfill
  {\color{White}\fontsize{9.6}{11.5}\selectfont Global architecture, with Australia as the worked example\par}
  \vspace{0.12cm}
  {\color{Mist}\fontsize{7.8}{9.5}\selectfont Integrated briefing current to \insertdate\ \textperiodcentered\ approx.\ 90 minutes}
  \note{[Timing] Framing block: 12 minutes including this slide. [Sources] Companion to The Sovereign Ledger integrated briefing, 5 August 2026.}
\end{frame}
\endgroup

% ============================ FRAMING ============================
\begin{frame}{The question is not who moves the last dollar}
  \claim{Who controls the path that determines whether the dollar can reach the final ledger?}
  \vspace{0.22cm}
  \begin{columns}[T,onlytextwidth]
    \column{0.48\textwidth}
      \rust{Formal authority}
      \begin{itemize}\tightitems
        \item Central bank issues the settlement asset
        \item Domestic law defines finality
        \item RTGS records the final debit and credit
      \end{itemize}
    \column{0.48\textwidth}
      \blue{Operational authority}
      \begin{itemize}\tightitems
        \item Identity and admission
        \item Routing, sequencing and liquidity locks
        \item Code, keys, data, recovery and exit
      \end{itemize}
  \end{columns}
  \source{RBA, \href{https://www.rba.gov.au/payments-and-infrastructure/rits/about.html}{RITS overview}; Bank of England, \href{https://www.bankofengland.co.uk/payment-and-settlement/a-brief-introduction-to-the-real-time-gross-settlement-system-and-chaps}{RTGS overview}.}
  \note{[Sources] RBA RITS overview; Bank of England RTGS overview.}
\end{frame}

\begin{frame}{Orders of magnitude} % FAST
  \begin{itemize}\tightitems
    \item \rust{Swift} moves value equivalent to world GDP every two to three days.
    \item \rust{CLS} handles FX settlement values on the order of US\$7 trillion on an average day.
    \item \rust{LCH SwapClear} clears about 95\% of vanilla OTC interest-rate swaps globally, in 28 currencies.
    \item \rust{Indeval's DAL\'I} settles over US\$250 billion of Mexican securities daily.
    \item \rust{Kinexys} (J.P.\ Morgan) reports over US\$4 trillion processed since inception.
  \end{itemize}
  \vspace{0.22cm}
  \begin{statement}
    These layers govern flow, not stock. Small design choices compound across these volumes daily.
  \end{statement}
  \source{Swift, July 2026 ledger announcement; CLS; LSEG SwapClear; Indeval/INFORMS; J.P.\ Morgan Kinexys.}
  \note{[Sources] Swift 9 July 2026 release (world-GDP claim); CLS published averages; LSEG SwapClear service page; Indeval Edelman materials; J.P. Morgan Kinexys newsroom.}
\end{frame}

\begin{frame}{Control of a settlement system is an eight-layer capability}
  \begin{tabularx}{\textwidth}{L{0.48cm}Y L{0.48cm}Y}
    \textcolor{Rust}{\textbf{1}} & Law and finality & \textcolor{Rust}{\textbf{5}} & Code, keys and specialist knowledge \\
    \addlinespace[0.35em]
    \textcolor{Rust}{\textbf{2}} & Rules, access and pricing & \textcolor{Rust}{\textbf{6}} & Connectivity, identity and certificates \\
    \addlinespace[0.35em]
    \textcolor{Rust}{\textbf{3}} & Authoritative money or title record & \textcolor{Rust}{\textbf{7}} & Data, routing and optimisation \\
    \addlinespace[0.35em]
    \textcolor{Rust}{\textbf{4}} & Live operations and recovery & \textcolor{Rust}{\textbf{8}} & Credible exit and substitution \\
  \end{tabularx}
  \vspace{0.30cm}
  \begin{statement}
    A jurisdiction can retain layers 1--3 while conceding layers 4--8 to external vendors and international utilities.
  \end{statement}
  \vspace{0.12cm}
  \muted{\fontsize{9.2}{11}\selectfont Layer tags \lyr{1}--\lyr{8} recur on the dependency and jurisdiction tables that follow.}
  \note{[Sources] Analytical framework developed in The Sovereign Ledger. The layer tags are reused on the dependency and jurisdiction tables so the framework does analytic work rather than decorating the introduction.}
\end{frame}

\begin{frame}{One payment contains several control events} % FAST
  \vspace{0.18cm}
  \begin{center}
    \fontsize{12}{14.5}\selectfont
    Initiate \quad $\rightarrow$ \quad authenticate \quad $\rightarrow$ \quad route \\
    \vspace{0.25cm}
    $\rightarrow$ \quad clear / optimise \quad $\rightarrow$ \quad reserve liquidity \\
    \vspace{0.25cm}
    $\rightarrow$ \quad \rust{settle in central-bank money}
  \end{center}
  \vspace{0.36cm}
  \begin{statement}
    The final RTGS entry can remain publicly operated even when every economically important choice was made upstream.
  \end{statement}
  \source{BIS CPMI, payment-system concepts; RBA RITS; Swift shared-ledger publications.}
  \note{[Sources] BIS CPMI conceptual framework; RBA RITS; Swift shared-ledger announcements.}
\end{frame}

\begin{frame}{Tokenisation adds records without removing law}
  \begin{columns}[T,onlytextwidth]
    \column{0.47\textwidth}
      \rust{Four records are called ``the ledger''}
      \begin{itemize}\tightitems
        \item Central-bank money ledger
        \item Securities or fund ownership record
        \item Commercial-bank liability ledger
        \item Coordination / orchestration ledger
      \end{itemize}
    \column{0.49\textwidth}
      \blue{The decisive questions}
      \begin{itemize}\tightitems
        \item Which record is legally authoritative?
        \item Who can upgrade or pause it?
        \item What happens in insolvency?
        \item Can the market operate without it?
      \end{itemize}
  \end{columns}
  \source{BIS Project Agor\'a final report; central-bank and FMI rulebooks surveyed in the companion report.}
  \note{[Sources] BIS Project Agorá final report; RBA, BoE, ECB, BOJ, SNB and FMI materials.}
\end{frame}

\begin{frame}{The central risk is a private constitution}
  \claim{A shared technical layer can become the place where market access, liquidity and permissible behaviour are actually governed.}
  \vspace{0.28cm}
  \begin{itemize}\tightitems
    \item No conspiracy is required: institutions rationally join every credible network.
    \item Network effects make the chosen identity, token and messaging standards difficult to escape.
    \item The result can resemble coordination because the same firms occupy every architecture.
  \end{itemize}
  \source{Ownership and participation analysis in The Sovereign Ledger; public corporate and FMI disclosures.}
  \note{[Sources] Companion report ownership atlas and participant mapping. This is an analytical inference, not an allegation of conspiracy.}
\end{frame}

\begin{frame}{Roadmap, method and a disclaimer}
  \begin{tabularx}{\textwidth}{L{0.62cm}Y L{0.62cm}Y}
    \textcolor{Rust}{\textbf{I}} & What the core already looks like & \textcolor{Rust}{\textbf{V}} & Swift and competing networks \\
    \addlinespace[0.28em]
    \textcolor{Rust}{\textbf{II}} & Australia's current position & \textcolor{Rust}{\textbf{VI}} & Technology and company lineage \\
    \addlinespace[0.28em]
    \textcolor{Rust}{\textbf{III}} & Eight Agor\'a jurisdictions & \textcolor{Rust}{\textbf{VII}} & ASX: failure and consequence \\
    \addlinespace[0.28em]
    \textcolor{Rust}{\textbf{IV}} & Agor\'a and its lineage & \textcolor{Rust}{\textbf{VIII}} & The test for any jurisdiction \\
  \end{tabularx}
  \vspace{0.22cm}
  \muted{\fontsize{9.4}{11.2}\selectfont Method: every assessment is tagged \tdoc, \tinf{} or \tunp.}
  \vspace{0.16cm}
  \begin{statement}
    This is not an argument for autarky: shared infrastructure is desirable --- on inspectable, substitutable, contestable terms.
  \end{statement}
  \note{[Timing] End of framing block at ~12 minutes. [Sources] Evidence-tier framework from The Sovereign Ledger companion report.}
\end{frame}

% ============================ PART I ============================
\sectionframe{PART I}{What the core already looks like}{10 minutes}

\begin{frame}{FMI ownership follows four recurring models}
  \begin{tabularx}{\textwidth}{L{2.55cm}Y Y}
    \toprule
    \textbf{Model} & \textbf{Examples} & \textbf{Primary capture risk} \\
    \midrule
    Public operator & RITS, Fedwire, RT2 & procurement and capability dependence \\
    \addlinespace[0.25em]
    User-owned utility & Swift, CLS, SIX, TCH & incumbent admission and rulebook power \\
    \addlinespace[0.25em]
    Listed FMI group & ASX, LSEG, Deutsche B\"orse, TMX & shareholder returns over resilience \\
    \addlinespace[0.25em]
    Venture-backed platform & Digital Asset, R3, Fireblocks & vendor, investor and exit concentration \\
    \bottomrule
  \end{tabularx}
  \source{Operator governance and shareholder disclosures cited in the companion report.}
  \note{[Sources] Swift, CLS, SIX, TCH, ASX, LSEG, Deutsche Börse, TMX, Digital Asset and R3 disclosures.}
\end{frame}

\begin{frame}{Securities control is not exchange ownership}
  \begin{columns}[T,onlytextwidth]
    \column{0.48\textwidth}
      \rust{Trading venue}
      \begin{itemize}\tightitems
        \item Matches buyers and sellers
        \item Can be competitively duplicated
      \end{itemize}
      \vspace{0.25cm}
      \rust{CCP}
      \begin{itemize}\tightitems
        \item Novates and nets obligations
        \item Controls margin and default management
      \end{itemize}
    \column{0.48\textwidth}
      \blue{CSD / registry}
      \begin{itemize}\tightitems
        \item Establishes or supports legal title
        \item Coordinates securities settlement
        \item Much harder to substitute in crisis
      \end{itemize}
  \end{columns}
  \source{CPMI-IOSCO PFMI; ASX CHESS; Euroclear CREST materials.}
  \note{[Sources] CPMI-IOSCO Principles for FMIs; ASX and Euroclear descriptions.}
\end{frame}

\begin{frame}{Commercial-bank money already settles at the core}
  \begin{itemize}\tightitems
    \item \rust{Euroclear Bank}, the world's largest international CSD, settles the cash leg in commercial-bank money: claims on Euroclear Bank itself.
    \item \rust{CLS}: funding and defunding move through central-bank RTGS accounts, but settlement occurs across accounts on the books of CLS Bank; CLSNet nets outside central-bank money.
    \item \rust{Correspondent banking} discharges cross-border obligations in commercial-bank money by construction.
  \end{itemize}
  \vspace{0.12cm}
  \begin{statement}
    The pure central-bank-money story was already porous. Tokenisation does not introduce commercial-bank settlement money; it makes it programmable. \hfill\tdoc
  \end{statement}
  \source{Euroclear Bank service descriptions; CLS settlement and CLSNet materials; BIS CPMI correspondent-banking reports.}
  \note{[Sources] Euroclear Bank (ICSD settlement in Euroclear Bank money); CLS settlement mechanics; CPMI. This slide answers the objection that commercial-bank settlement money is novel.}
\end{frame}

\begin{frame}{Message-layer control is demonstrated power}
  \begin{itemize}\tightitems
    \item 2022: EU sanctions disconnected major Russian banks from Swift --- the messaging layer used, openly, as a policy instrument.
    \item CIPS (2015) and the e-CNY are the most developed state-built responses; mBridge (Part IV) is the multi-central-bank variant.
    \item Every jurisdiction absorbed the same lesson: whoever operates the coordination layer can exclude.
  \end{itemize}
  \vspace{0.28cm}
  \begin{statement}
    No further proof is required that the layers above final settlement carry decisive power. The only open question is who holds them next. \hfill\tdoc
  \end{statement}
  \source{EU Council sanctions decisions (2022); PBoC CIPS materials; BIS mBridge documentation.}
  \note{[Sources] EU Council Regulation amendments removing named Russian banks from Swift-provided messaging; CIPS; PBoC. Framed as a design lesson shared by all jurisdictions, not a bloc-specific grievance.}
\end{frame}

\begin{frame}{The physical layer concentrates as well}
  \begin{itemize}\tightitems
    \item FMIs and their vendors increasingly run on a small set of hyperscalers and shared software supply chains --- layers \lyr{4} and \lyr{6} in practice.
    \item 19 July 2024: a single faulty CrowdStrike update simultaneously disrupted airlines, banks and exchanges worldwide --- one vendor, global propagation.
    \item Cloud regions, HSM fleets, certificate authorities and privileged tooling sit outside most jurisdictions' practical reach.
  \end{itemize}
  \vspace{0.28cm}
  \begin{statement}
    Layer analysis must include hosting, keys and endpoint software --- not only rulebooks. \hfill\tdoc
  \end{statement}
  \source{Public incident reporting, July 2024; FSB third-party dependency workstream.}
  \note{[Sources] Widely documented 19 July 2024 CrowdStrike outage; FSB and national-regulator third-party risk publications.}
\end{frame}

% ============================ PART II ============================
\sectionframe{PART II}{Australia's current position}{10 minutes}

\begin{frame}{Australia owns the monetary core}
  \begin{columns}[T,onlytextwidth]
    \column{0.52\textwidth}
      \begin{itemize}\tightitems
        \item RBA owns and operates \rust{RITS}
        \item ESA transfers provide final wholesale AUD settlement
        \item FSS supplies 24/7 central-bank settlement for NPP
        \item CHESS remains the legal-title and equity-settlement spine
      \end{itemize}
    \column{0.44\textwidth}
      \begin{statement}
        Australia issues, records and finally settles its own currency. That does not by itself confer end-to-end operational control.
      \end{statement}
  \end{columns}
  \source{RBA, \href{https://www.rba.gov.au/payments-and-infrastructure/rits/about.html}{RITS}; RBA, \href{https://www.rba.gov.au/publications/bulletin/2018/sep/the-new-payments-platform-and-fast-settlement-service.html}{FSS}; ASX, CHESS settlement.}
  \note{[Sources] RBA RITS and FSS publications; ASX CHESS settlement description.}
\end{frame}

\begin{frame}{Australia depends above the final ledger}
  \begin{tabularx}{\textwidth}{L{2.25cm}Y Y L{0.95cm}}
    \toprule
    \textbf{Layer} & \textbf{Important provider} & \textbf{Dependency} & \textbf{Tags} \\
    \midrule
    High-value messaging & Swift & standards, connectivity, certificates & \lyr{6} \\
    Retail fast payments & AP+ / Swift / FSS & scheme, gateway and settlement integration & \lyr{2}\,\lyr{6} \\
    COIN connectivity & TNS & network operations and support & \lyr{6} \\
    Equity replacement & TCS BaNCS & foreign product and specialist capability & \lyr{5} \\
    Future orchestration & undecided & routing, locks, data and workflow authority & \lyr{2}\,\lyr{7} \\
    \bottomrule
  \end{tabularx}
  \source{RBA, AusPayNet, AP+, Swift, TNS, ASX and TCS disclosures.}
  \note{[Sources] RBA; AusPayNet COIN; TNS; ASX CHESS replacement; TCS BaNCS. Layer tags refer to the eight-layer framework.}
\end{frame}

\begin{frame}{The largest concession already happened}
  \begin{itemize}\tightitems
    \item The large majority of AUD OTC interest-rate derivatives clears at \rust{LCH SwapClear} in London ($\sim$95\% of vanilla OTC swaps globally).
    \item The CFR's 2012--2015 location-policy debate settled on overseas clearing under safeguards (licensing, RBA standards, Bank of England cooperation).
    \item ASX Clear (Futures) offers a domestic AUD/NZD OTC service; the global liquidity pool sits offshore.
  \end{itemize}
  \vspace{0.10cm}
  \begin{statement}
    Margin, default management and crisis decisions for the core AUD rates market already sit with a London CCP. The tokenisation debate starts here, not from autonomy. \hfill\tdoc
  \end{statement}
  \source{CFR, \href{https://www.cfr.gov.au/publications/policy-statements-and-other-reports/2013/report-on-the-australian-otc-derivatives-market-july/}{OTC derivatives reports}; RBA Bulletin (2018); LSEG, \href{https://www.lseg.com/en/post-trade/clearing/lch-services/swapclear}{SwapClear}; ASX OTC clearing.}
  \note{[Sources] CFR OTC derivatives market reports; RBA 2018 Bulletin on trade-repository data; LSEG SwapClear; ASX OTC clearing service page. Supports the second-edition thesis: tokenisation reveals existing external dependence rather than creating it.}
\end{frame}

\begin{frame}{RITS modernisation is the live decision}
  \begin{columns}[T,onlytextwidth]
    \column{0.48\textwidth}
      \rust{Opportunity}
      \begin{itemize}\tightitems
        \item Modern interfaces and data
        \item Longer operating hours
        \item More settlement in central-bank money
        \item Multiple synchronisers, not one franchise
      \end{itemize}
    \column{0.48\textwidth}
      \blue{Present weakness}
      \begin{itemize}\tightitems
        \item Three PFMI principles partly observed
        \item Core Modernisation delays
        \item Congested change portfolio
        \item Scarce specialist capacity
      \end{itemize}
  \end{columns}
  \source{RBA, \href{https://www.rba.gov.au/payments-and-infrastructure/rits/self-assessments/2026/ratings-and-recommendations.html}{2026 ratings} and \href{https://www.rba.gov.au/payments-and-infrastructure/rits/self-assessments/2026/material-developments.html}{material developments}.}
  \note{[Sources] RBA 2026 RITS assessment, ratings and material developments.}
\end{frame}

% ============================ PART III ============================
\sectionframe{PART III}{Eight Agor\'a jurisdictions}{11 minutes}

\begin{frame}{The euro area: public anchor, private machinery}
  \begin{columns}[T,onlytextwidth]
    \column{0.54\textwidth}
      \rust{Existing stack}
      \begin{itemize}\tightitems
        \item T2 / T2S / TIPS central-bank services
        \item Euroclear and Clearstream post-trade groups
        \item Eurex and LCH clearing power
      \end{itemize}
      \blue{Tokenisation route}
      \begin{itemize}\tightitems
        \item 2024 trials: Bundesbank Trigger, Banque de France DL3S, TIPS Hash-Link
        \item \rust{Pontes} bridges market DLTs to TARGET; \rust{Appia} is the integrated long run
      \end{itemize}
    \column{0.42\textwidth}
      \begin{statement}
        The euro's monetary anchor is public; much of its collateral and securities machinery is not.
      \end{statement}
  \end{columns}
  \source{ECB TARGET, \href{https://www.ecb.europa.eu/paym/target/pontes/html/index.en.html}{Pontes} and Appia; Euroclear and Deutsche B\"orse disclosures.}
  \note{[Sources] ECB exploratory-work report, Pontes and Appia roadmap; Euroclear; Deutsche Börse/Clearstream. Merges the two euro slides from the earlier edition.}
\end{frame}

\begin{frame}{Japan separates bank money from orchestration}
  \begin{itemize}\tightitems
    \item BOJ-NET settles central-bank money and JGBs; Zengin-Net clears retail/interbank payments.
    \item DCJPY is live with one issuing bank: GMO Aozora.
    \item DCP transmits instructions; the issuing bank's Financial Zone record is characterised as the deposit.
    \item Multi-bank interbank settlement remained an FSA-supported proof of concept in April 2026.
  \end{itemize}
  \begin{statement}
    Public evidence does not establish Article 73 insolvency protection for DCP-level netting. \hfill\tinf
  \end{statement}
  \source{BOJ; DCP legal/product disclosures; Japan Payment Services Act arts 64, 71, 73.}
  \note{[Sources] BOJ; DCP about/eps/company pages and April 2026 PoC; Japan Payment Services Act. Qualification: unpublished agreements could route positions into a protected mechanism.}
\end{frame}

\begin{frame}{Korea is testing tokenised deposits under a public core} % FAST
  \begin{itemize}\tightitems
    \item BOK-Wire+ remains the high-value central-bank-money system.
    \item The Bank of Korea has explored wholesale CBDC and deposit-token architectures.
    \item Six Korean banks in Agor\'a provide broad domestic coverage: Hana, IBK, KB, NongHyup, Shinhan and Woori.
    \item The policy question is whether a shared platform complements BOK-Wire+ or becomes the indispensable access layer above it.
  \end{itemize}
  \source{Bank of Korea speeches and Agor\'a participant package.}
  \note{[Sources] Bank of Korea Governor speeches; BIS/IIF Agorá participant list.}
\end{frame}

\begin{frame}{Mexico tests cross-border reach without surrendering SPEI}
  \begin{itemize}\tightitems
    \item Banco de M\'exico operates SPEI and the account system SIAC; SPID handles USD payments between legal entities.
    \item \rust{Indeval}, the private CSD, operates the DAL\'I securities-settlement system, linked to Banco de M\'exico's systems.
    \item Banorte supplies a domestic bank perspective; Monex adds specialist FX and direct CLS participation.
    \item BBVA connects the peso domain to a global group and published Agor\'a validation work.
  \end{itemize}
  \begin{statement}
    The test is whether domestic peso discharge remains anchored to Banco de M\'exico money.
  \end{statement}
  \source{Banco de M\'exico system descriptions; \href{https://www.indeval.com.mx/en/product?p=clearing-settlement}{Indeval, DAL\'I}; BBVA; BIS/IIF Agor\'a participant list.}
  \note{[Sources] Banco de México (SPEI, SIAC, SPID); Indeval (DALÍ operator); BBVA Agorá publication; BIS/IIF list. Corrects the earlier edition, which attributed DALÍ's operation to Banco de México.}
\end{frame}

\begin{frame}{Switzerland operates both public and private models}
  \begin{itemize}\tightitems
    \item SIC is operated by SIX under the SNB mandate.
    \item SNB wholesale CBDC has settled live transactions on SDX since 2023.
    \item Helvetia now compares integrated wCBDC with a production RTGS link.
    \item Fnality-like private settlement and multi-bank deposit tokens remain legally distinct from SNB money.
  \end{itemize}
  \begin{statement}
    Switzerland is testing architecture without pretending every Swiss-franc token is the same liability.
  \end{statement}
  \source{SNB Project Helvetia; SIX/SDX; Swiss multi-bank deposit-token PoC.}
  \note{[Sources] SNB Helvetia pages and extension release; SIX/SDX; Sygnum deposit-token PoC.}
\end{frame}

\begin{frame}{The UK renewed the core and made synchronisation a role}
  \begin{itemize}\tightitems
    \item The Bank of England's \rust{RT2} went live in April 2025.
    \item Mastercard's VocaLink runs key retail infrastructure; Euroclear UK \& International operates CREST.
    \item The Synchronisation Lab formalises external operators coordinating RTGS with asset ledgers.
    \item Fnality supplies reserve-backed settlement tokens; Lloyds and NatWest link deposit-token work to Agor\'a.
  \end{itemize}
  \begin{statement}
    A synchronisation operator need not own RTGS to control when an instruction is released. Its power lies in conditions, timing and earmarks.
  \end{statement}
  \source{Bank of England RT2 and Synchronisation Lab; Mastercard/VocaLink; Euroclear EUI; Fnality; Lloyds and NatWest materials.}
  \note{[Sources] Bank of England RT2 launch and Synchronisation Lab; Mastercard acquisition; Euroclear EUI annual report; Fnality; Lloyds tokenised deposit/gilt transaction. Merges the two UK slides from the earlier edition.}
\end{frame}

\begin{frame}{The United States privatises dollar governance}
  \begin{itemize}\tightitems
    \item Fedwire and FedNow provide public settlement rails.
    \item TCH owns CHIPS, RTP and EPN on behalf of the large US commercial banks that own it.
    \item DTCC controls the central US securities post-trade chain.
    \item Kinexys, Citi and BNY operate proprietary commercial-bank-money ledgers.
  \end{itemize}
  \begin{statement}
    The dollar's monetary anchor is secure; access to programmable dollars may still be governed through private tollbooths.
  \end{statement}
  \source{Federal Reserve; TCH; DTCC; JPMorgan, Citi and BNY token-service disclosures.}
  \note{[Sources] Federal Reserve; TCH (owner count varies with membership --- verify against the current shareholder page before quoting a number); DTCC; JPMorgan/Citi/BNY.}
\end{frame}

\begin{frame}{Canada combines a public core with imported optimisation}
  \begin{itemize}\tightitems
    \item Payments Canada owns Lynx; settlement occurs in Bank of Canada money.
    \item TMX owns CDS and CDCC.
    \item TMX and Clearstream jointly operate the Canadian Collateral Management Service.
    \item Canada joined Agor\'a in May 2026, after the seven-currency prototype.
  \end{itemize}
  \begin{statement}
    CCMS shows how a foreign-controlled optimiser can shape domestic repo without owning the central bank or CSD.
  \end{statement}
  \source{Payments Canada; Bank of Canada Agor\'a announcement; TMX/Clearstream CCMS.}
  \note{[Sources] Payments Canada Lynx; Bank of Canada (May 2026 accession alongside the 27 May findings report); TMX/Clearstream CCMS launch.}
\end{frame}

\begin{frame}{Eight jurisdictions, eight constitutional positions}
  \fontsize{8.0}{9.6}\selectfont
  \begin{tabularx}{\textwidth}{L{1.28cm}L{1.95cm}Y Y L{1.12cm}}
    \toprule
    \textbf{Area} & \textbf{Public anchor} & \textbf{Tokenisation path} & \textbf{Main exposure} & \textbf{Tags} \\
    \midrule
    Euro & TARGET & Pontes / Appia & private post-trade concentration & \lyr{3}\,\lyr{7} \\
    Japan & BOJ-NET & DCJPY / sandbox & multi-bank discharge & \lyr{1}\,\lyr{3} \\
    Korea & BOK-Wire+ & CBDC/deposit tests & access-layer dependence & \lyr{2}\,\lyr{8} \\
    Mexico & SPEI & Agor\'a / bank tokens & peso settlement anchor & \lyr{1}\,\lyr{3} \\
    Switzerland & SIC / SNB wCBDC & Helvetia / SDX & SIX and platform control & \lyr{2}\,\lyr{4} \\
    UK & RT2 & sync / Fnality & external operator and earmarks & \lyr{4}\,\lyr{7} \\
    US & Fedwire & TCH / Kinexys / DTCC & private dollar governance & \lyr{2} \\
    Canada & Lynx & Agor\'a / CCMS & imported collateral optimisation & \lyr{7} \\
    \bottomrule
  \end{tabularx}
  \source{Central-bank and FMI sources summarised on preceding slides.}
  \note{[Sources] Preceding jurisdiction slides and companion report. Layer tags map each exposure back to the eight-layer framework.}
\end{frame}

% ============================ PART IV ============================
\sectionframe{PART IV}{Agor\'a and its experimental lineage}{12 minutes}

\begin{frame}{Agor\'a joins bank deposits to reserve domains}
  \begin{center}
    \fontsize{13}{15.5}\selectfont
    Tokenised commercial-bank deposits \\
    \vspace{0.18cm}
    $\downarrow$ \quad shared payment contracts and atomic conditions \quad $\downarrow$ \\
    \vspace{0.18cm}
    Separate currency-specific central-bank reserve domains
  \end{center}
  \vspace{0.32cm}
  \begin{itemize}\tightitems
    \item Seven original currency areas; Canada joined in May 2026 after the completed prototype.
    \item Besu-based experimental architecture with ISO 20022 and legal-entity identity.
    \item July 2026 real-value phase: 17 scenarios, about CHF800,000 aggregate value, roughly 80 seconds average initiation-to-settlement.
  \end{itemize}
  \source{BIS, \href{https://www.bis.org/publ/othp110.pdf}{Agor\'a final report} (27 May 2026); BIS, \href{https://www.bis.org/about/bisih/topics/fmis/agora/rvt.htm}{real-value testing summary}.}
  \note{[Sources] BIS Project Agorá findings report of 27 May 2026 and July 2026 real-value testing summary (published 30 July 2026).}
\end{frame}

\begin{frame}{Agor\'a proved mechanics, not a constitution}
  \begin{columns}[T,onlytextwidth]
    \column{0.47\textwidth}
      \rust{What it demonstrated}
      \begin{itemize}\tightitems
        \item Atomic multi-currency transactions
        \item Shared programmable workflows
        \item Technical operation across institutions
        \item Real-value test execution
      \end{itemize}
    \column{0.49\textwidth}
      \blue{What remains open}
      \begin{itemize}\tightitems
        \item Production owner and operator
        \item Rulebook and legal designation
        \item Loss allocation and insolvency
        \item Validator control, recovery and exit
      \end{itemize}
  \end{columns}
  \source{BIS Agor\'a final report and real-value testing documentation.}
  \note{[Sources] BIS Project Agorá final report; BIS real-value testing page.}
\end{frame}

\begin{frame}{The common ledger does not erase RTGS}
  \begin{itemize}\tightitems
    \item Reserve tokens must be funded from, and redeemable into, the relevant central-bank accounts.
    \item Commercial-bank tokens remain claims on named banks.
    \item Cross-currency atomicity can occur on the shared platform while domestic finality remains anchored to each jurisdiction's money.
    \item A production design could nevertheless make the common layer operationally indispensable.
  \end{itemize}
  \begin{statement}
    A currency authority keeps control only if it can isolate its domain and continue settlement without the shared layer.
  \end{statement}
  \source{BIS Agor\'a architecture and companion-report legal analysis.}
  \note{[Sources] BIS Project Agorá final report. Last sentence is the report's control test.}
\end{frame}

\begin{frame}{Forty-two firms shaped the prototype}
  \begin{columns}[T,onlytextwidth]
    \column{0.48\textwidth}
      \rust{Concentrated capabilities}
      \begin{itemize}\tightitems
        \item Global correspondents and FX dealers
        \item Japanese and Korean domestic bank clusters
        \item Custodians and tokenised-deposit operators
        \item Eurex, Euroclear and SIX Digital Exchange
        \item Visa, Mastercard and Swift
      \end{itemize}
    \column{0.48\textwidth}
      \blue{Under-represented interests}
      \begin{itemize}\tightitems
        \item End users and corporates
        \item Smaller-bank associations
        \item Independent competition voices
        \item Public-interest and civil-society representatives
      \end{itemize}
  \end{columns}
  \source{BIS/IIF private-sector participation package and final participant list.}
  \note{[Sources] BIS/IIF Agorá application package; final report participant list.}
\end{frame}

\begin{frame}{Participation is not ownership}
  \claim{The 42 firms were a design and testing coalition---not shareholders in a disclosed production utility.}
  \vspace{0.25cm}
  \begin{itemize}\tightitems
    \item Real-value testing was conducted by 28 institutions on behalf of all participants, in transactions of CHF 9,000--125,000.
    \item Project status therefore proves workstream access, not execution of a real-value transaction --- and not endorsement of every design conclusion.
    \item Only a few specialised roles were uniquely disclosed: Eurex as sole CCP; SIX and Kaleido in later test roles.
    \item The future operator, ownership and governance remain undecided.
  \end{itemize}
  \source{BIS/IIF application package; BIS Agor\'a final and \href{https://www.bis.org/about/bisih/topics/fmis/agora/rvt.htm}{real-value testing} reports.}
  \note{[Sources] BIS/IIF participation package; BIS Project Agorá reports; RVT summary (28 institutions, CHF 9k–125k values).}
\end{frame}

\begin{frame}{Dunbar tested one multi-CBDC platform} % FAST
  \begin{itemize}\tightitems
    \item RBA, Bank Negara Malaysia, MAS and South African Reserve Bank.
    \item Two prototypes: \rust{Corda/R3} and \blue{Quorum/Partior}.
    \item Direct cross-border transactions in several hypothetical CBDCs.
    \item Governance, non-resident access and central-bank autonomy were the hard problems.
  \end{itemize}
  \begin{statement}
    Dunbar asked whether several national monies could inhabit one platform without a platform super-user overruling the issuers.
  \end{statement}
  \source{BIS, \href{https://www.bis.org/publ/othp47.htm}{Project Dunbar final report}.}
  \note{[Sources] BIS Project Dunbar final report.}
\end{frame}

\begin{frame}{Jura put real central-bank money on Corda} % FAST
  \begin{itemize}\tightitems
    \item Banque de France, SNB and BIS with SIX, UBS, Credit Suisse, Natixis, R3 and Accenture.
    \item Real-value euro and Swiss-franc wCBDC plus tokenised commercial paper.
    \item Subnetworks and dual-notary signing preserved currency-specific vetoes.
  \end{itemize}
  \begin{statement}
    A central bank can retain issuance authority while depending on a third-party platform for availability and software evolution.
  \end{statement}
  \source{BIS, \href{https://www.bis.org/publ/othp44.htm}{Project Jura report}.}
  \note{[Sources] BIS Project Jura report.}
\end{frame}

\begin{frame}{mBridge is the central-bank-built alternative}
  \begin{itemize}\tightitems
    \item Hong Kong, Thailand, China, UAE and later Saudi Arabia.
    \item Custom mBridge Ledger with a validating node for each founding authority.
    \item Real-value bank participation; MVP reached in 2024; BIS handed the project to participating authorities in October 2024.
    \item The handover followed public debate about whether multi-CBDC platforms could route around established payment controls. \tinf
  \end{itemize}
  \begin{statement}
    mBridge can reduce dependence on correspondent chains and Swift---but validator control does not resolve sanctions, admission or rulebook politics.
  \end{statement}
  \source{BIS, \href{https://www.bis.org/about/bisih/topics/cbdc/mcbdc_bridge.htm}{mBridge project page} and 2022 report.}
  \note{[Sources] BIS mBridge project page, 2022 report and 2024 handover announcement; contemporaneous public commentary on the exit (tagged inference: motives were not fully disclosed).}
\end{frame}

\begin{frame}{Three constitutional models are competing}
  \fontsize{9.6}{11.4}\selectfont
  \begin{tabularx}{\textwidth}{L{2.05cm}Y Y}
    \toprule
    \textbf{Model} & \textbf{Examples} & \textbf{Concentration danger} \\
    \midrule
    Common multi-issuer platform & Dunbar, mBridge & common governance and validator layer \\
    Shared bank/public workflow & Agor\'a, Swift ledger & incumbent orchestration and identity \\
    Interlinked national ledgers & Cedar--Ubin+, Pontes, Meridian FX, Helvetia & bridge, oracle and synchroniser power \\
    \bottomrule
  \end{tabularx}
  \begin{statement}
    There is no inevitable one-ledger future. The contest is where coordination authority sits.
  \end{statement}
  \source{BIS Innovation Hub project reports; ECB Pontes; SNB Helvetia; New York Fed Cedar. Project-by-project control questions in the appendix.}
  \note{[Sources] Comparative synthesis of BIS, central-bank and Swift project documentation. The per-project mechanism table (Mariana, Cedar–Ubin+, Helvetia, Pontes, Meridian FX) moved to the appendix.}
\end{frame}

% ============================ PART V ============================
\sectionframe{PART V}{Swift and competing networks}{12 minutes}

\begin{frame}{Swift is moving from messages to commitments}
  \begin{itemize}\tightitems
    \item September 2025: shared-ledger plan designed with more than 30 institutions.
    \item March 2026: Besu-based MVP records and validates interbank payment commitments.
    \item Swift operates the ledger; banks retain keys, deposit assets, compliance and ultimate funding.
    \item July 2026: controlled initial-use phase announced with 17 banks.
  \end{itemize}
  \begin{statement}
    Swift is not becoming the settlement asset. It is trying to become the trusted workflow through which settlement is arranged.
  \end{statement}
  \source{Swift, shared-ledger announcements of September 2025, March 2026 and 9 July 2026.}
  \note{[Sources] Swift shared-ledger announcement; MVP architecture; 17-bank initial-use cohort.}
\end{frame}

\begin{frame}{The 17-bank cohort is globally distributed} % FAST
  \fontsize{10.3}{12.4}\selectfont
  \begin{columns}[T,onlytextwidth]
    \column{0.32\textwidth}
      \rust{Europe / Americas}
      \begin{itemize}\verytightitems
        \item BNP Paribas
        \item BNY
        \item Citi
        \item HSBC
        \item Ita\'u
        \item Lloyds
        \item UBS
        \item Wells Fargo
      \end{itemize}
    \column{0.32\textwidth}
      \blue{Asia-Pacific}
      \begin{itemize}\verytightitems
        \item ANZ
        \item DBS
        \item MUFG
        \item OCBC
        \item UOB
      \end{itemize}
    \column{0.32\textwidth}
      \rust{Middle East / Africa}
      \begin{itemize}\verytightitems
        \item First Abu Dhabi Bank
        \item FirstRand
        \item Mashreq
        \item Standard Chartered
      \end{itemize}
  \end{columns}
  \source{Swift, \href{https://www.swift.com/news-events/press-releases/swifts-blockchain-ledger-ready-use-17-banks-set-pioneer-tokenised-cross-border-payments-trusted-global-infrastructure}{9 July 2026 initial-use announcement}.}
  \note{[Sources] Swift 9 July 2026 press release.}
\end{frame}

\begin{frame}{JPMorgan is absent from 17, not from the strategy}
  \begin{itemize}\tightitems
    \item JPMorgan was in Swift's earlier design group, but not the July 2026 cohort of 17.
    \item It is an Agor\'a participant and owns \rust{Kinexys}.
    \item Kinexys operates payments, tokenised assets and information networks.
    \item JPM Coin has operated on proprietary infrastructure and Base; native Canton issuance is planned.
  \end{itemize}
  \begin{statement}
    JPMorgan is hedging architectures, not choosing one.
  \end{statement}
  \source{Swift design and cohort releases; JPMorgan Kinexys/JPM Coin; Canton deployment announcement.}
  \note{[Sources] Swift; JPMorgan Kinexys and JPM Coin; Canton/JPM Coin January 2026 announcement.}
\end{frame}

\begin{frame}{CLS is the mature comparator}
  \begin{itemize}\tightitems
    \item Bank-owned cross-border utility with powerful network effects.
    \item Payment-versus-payment settlement eliminates principal risk for eligible FX.
    \item Funding moves through each currency's central-bank RTGS account; settlement itself occurs across accounts on the books of CLS Bank.
    \item Swift provides connectivity; cooperation is public, not evidence of merger.
  \end{itemize}
  \begin{statement}
    CLS shows how a private utility can become indispensable without replacing national central-bank money.
  \end{statement}
  \source{CLSSettlement, shareholder and third-party materials; Swift/CLS connectivity.}
  \note{[Sources] CLS product, shareholder and group-structure pages; Swift and CLS page. The books-of-CLS-Bank mechanics echo the Part I commercial-bank-money slide.}
\end{frame}

\begin{frame}{TCH proposes bank-owned tokenised deposits}
  \begin{itemize}\tightitems
    \item June 2026 proposal backed by 17 named financial institutions.
    \item On-chain clearing and settlement of tokenised bank deposits.
    \item Intended connectivity to CHIPS and RTP.
    \item TCH is owned by its large US member banks.
  \end{itemize}
  \begin{statement}
    The bank branch of programmable dollars preserves deposit franchises and owner-bank governance.
  \end{statement}
  \source{The Clearing House, \href{https://www.theclearinghouse.org/payment-systems/Articles/2026/06/Major-Financial-Institutions-Unveil-Bank-Led-On-Chain-Money-Initiative}{5 June 2026 initiative}.}
  \note{[Sources] TCH bank-led on-chain money initiative. Owner count varies with membership; verify against the current shareholder page before quoting a figure.}
\end{frame}

\begin{frame}{Open USD proposes a distribution coalition}
  \begin{itemize}\tightitems
    \item Announced June 2026 by Open Standard / OnePay; more than 140 businesses named, not yet a live stablecoin at the cutoff.
    \item Fee-free uncapped mint and redemption proposed; reserve earnings shared with participating distributors after a management fee.
    \item Context: the GENIUS Act (July 2025) bars paying yield to stablecoin holders --- revenue-sharing with \emph{distributors} is the compliant substitute.
  \end{itemize}
  \begin{statement}
    Open USD competes by paying platforms for distribution; TCH competes by preserving bank deposits.
  \end{statement}
  \source{\href{https://joinopenstandard.com/}{Open Standard}; \href{https://www.onepay.com/newsroom/introducing-open-usd}{OnePay launch announcement}; GENIUS Act (2025).}
  \note{[Sources] Open Standard site and OnePay announcement; GENIUS Act statutory text on holder yield.}
\end{frame}

\begin{frame}{Tokenised MMFs become payment-adjacent collateral} % FAST
  \begin{itemize}\tightitems
    \item BUIDL, BENJI/FOBXX and USYC combine fund claims with transferable tokens.
    \item They can accelerate subscription, redemption and collateral mobility.
    \item They are securities/fund interests---not bank deposits or central-bank money.
    \item NAV, gates, official ownership records and eligible-investor restrictions remain decisive.
  \end{itemize}
  \begin{statement}
    A token can trade like cash operationally while retaining investment and insolvency risk legally.
  \end{statement}
  \source{Franklin Templeton BENJI; Circle USYC; Securitize/BlackRock BUIDL materials.}
  \note{[Sources] Franklin Templeton; Circle; Securitize/BlackRock.}
\end{frame}

\begin{frame}{ERC standards solve interfaces, not finality} % FAST
  \begin{tabularx}{\textwidth}{L{2.1cm}Y Y}
    \toprule
    \textbf{Standard} & \textbf{Function} & \textbf{Does not establish} \\
    \midrule
    ERC-20 & fungible token interface & legal title or redemption \\
    ERC-4626 & synchronous tokenised vaults & investor eligibility or NAV certainty \\
    ERC-7540 & asynchronous request lifecycle & settlement asset or insolvency priority \\
    ERC-7575 & multi-asset vault extensions & authoritative registry status \\
    \bottomrule
  \end{tabularx}
  \source{Ethereum Improvement Proposals 20, 4626, 7540 and 7575.}
  \note{[Sources] Official Ethereum Improvement Proposals.}
\end{frame}

\begin{frame}{Identity is the quiet gatekeeper}
  \fontsize{9.7}{11.6}\selectfont
  \begin{columns}[T,onlytextwidth]
    \column{0.48\textwidth}
      \rust{Architectures}
      \begin{itemize}\tightitems
        \item Central registry: one provider owns eligibility data
        \item Reusable issuer-signed credentials
        \item Federated multi-authority verification
        \item Zero-knowledge proofs of eligibility
      \end{itemize}
    \column{0.48\textwidth}
      \blue{Securitize, the live instantiation}
      \begin{itemize}\tightitems
        \item Transfer agent, broker-dealer/ATS, fund admin
        \item Allowlists encode jurisdiction and eligibility
        \item BlackRock: 2024 round and a board seat
        \item Listed as SECZ, July 2026 (Cantor-sponsored)
      \end{itemize}
  \end{columns}
  \vspace{0.08cm}
  \begin{statement}
    Whoever controls the identity-to-wallet mapping can become the private gatekeeper to tokenised markets.
  \end{statement}
  \source{W3C Verifiable Credentials 2.0; Securitize corporate and investor releases; SEC filings.}
  \note{[Sources] W3C VC Data Model 2.0; Securitize BlackRock-led round and July 2026 listing; SEC filings. Merges the earlier identity and Securitize slides. Companion point: portability matters only if credentials, revocation and recovery remain contestable.}
\end{frame}

% ============================ PART VI ============================
\sectionframe{PART VI}{Technology and company lineage}{11 minutes}

\begin{frame}{Technology choice is constitutional choice}
  \fontsize{9.1}{10.9}\selectfont
  \begin{tabularx}{\textwidth}{L{1.5cm}Y Y Y}
    \toprule
    & \textbf{Data model} & \textbf{Strongest feature} & \textbf{Capture point} \\
    \midrule
    Besu & shared EVM state & standards and labour pool & validators, privacy, upgrades \\
    Canton & party-specific contracts & privacy and cross-app atomicity & synchroniser, validators, Daml \\
    Corda & shared bilateral states & regulated workflows & notary and network manager \\
    Fabric & channels / private data & modular enterprise control & ordering and membership service \\
    Public chain & replicated state & reach and composability & validators, bridges, RPC, custody \\
    \bottomrule
  \end{tabularx}
  \source{LF Decentralized Trust; Digital Asset; R3; Ethereum documentation.}
  \note{[Sources] Besu, Fabric, Canton, Corda and Ethereum official documentation.}
\end{frame}

\begin{frame}{Besu is the institutional common denominator}
  \begin{itemize}\tightitems
    \item Apache-2.0 Ethereum client under LF Decentralized Trust.
    \item Public Ethereum or permissioned QBFT/IBFT networks.
    \item EVM compatibility exposes mature tools and token standards.
    \item Used differently by Agor\'a, Swift, Fnality, DTCC and JSCC.
  \end{itemize}
  \begin{statement}
    Open code reduces code lock-in; it does not eliminate control over validators, hosting, privacy services, keys or releases.
  \end{statement}
  \source{LF Decentralized Trust, \href{https://www.lfdecentralizedtrust.org/projects/besu}{Hyperledger Besu}; operator disclosures.}
  \note{[Sources] LF Decentralized Trust Besu; BIS Agorá; Swift; Fnality; DTCC; JSCC.}
\end{frame}

\begin{frame}{Fnality proves Besu can settle regulated value} % FAST
  \begin{itemize}\tightitems
    \item Sterling Fnality uses permissioned Ethereum/Besu.
    \item Tokens are fully backed through a Bank of England omnibus-account arrangement.
    \item The system is recognised and designated for settlement finality.
    \item Twenty-four banks/FMIs publicly described as backing the group.
  \end{itemize}
  \begin{statement}
    Production technology does not make the token central-bank money; each currency needs its own legal and central-bank arrangement.
  \end{statement}
  \source{Fnality; Bank of England RTGS/CHAPS annual report and recognition materials.}
  \note{[Sources] Fnality system descriptions; Bank of England treatment of Fnality.}
\end{frame}

\begin{frame}{Digital Asset came from the DRW orbit}
  \begin{center}
    \fontsize{11.5}{14}\selectfont
    DRW / Don Wilson \\
    \vspace{0.16cm}
    $\swarrow$ \hspace{2.2cm} $\searrow$ \\
    \vspace{0.12cm}
    Cumberland crypto liquidity \hspace{0.8cm} DRW investment activity \\
    \vspace{0.16cm}
    Eric Saraniecki \hspace{2.0cm} Yuval Rooz \\
    \vspace{0.23cm}
    $\searrow$ \hspace{2.0cm} $\swarrow$ \\
    \vspace{0.10cm}
    \rust{Digital Asset} \quad + \quad Sunil Hirani \quad + \quad Shaul Kfir
  \end{center}
  \vspace{0.25cm}
  \begin{itemize}\verytightitems
    \item Trading and crypto-market liquidity
    \item Derivatives workflow and clearing
    \item Research cryptography
  \end{itemize}
  \source{Digital Asset founder biographies and March 2015 launch; DRW, Don Wilson profile.}
  \note{[Sources] Digital Asset about page; 11 March 2015 launch release; DRW Don Wilson profile.}
\end{frame}

\begin{frame}{Founding record and the Masters era}
  \fontsize{10.2}{12.2}\selectfont
  \begin{columns}[T,onlytextwidth]
    \column{0.48\textwidth}
      \rust{March 2015 announcement}
      \begin{itemize}\tightitems
        \item Founders: Don Wilson and Sunil Hirani
        \item Technology described as their brainchild
        \item Both joined the board
      \end{itemize}
    \column{0.48\textwidth}
      \blue{Current corporate history}
      \begin{itemize}\tightitems
        \item Co-founders: Yuval Rooz, Eric Saraniecki, Shaul Kfir
        \item Masters: CEO from the March 2015 launch; not an original founder
        \item Reframed DLT as regulated-market infrastructure; led the ASX relationship
      \end{itemize}
  \end{columns}
  \vspace{0.08cm}
  \begin{statement}
    Best reconstruction: Wilson/Hirani financial founders; Rooz/Saraniecki/Kfir the operating-technical founding group; Masters supplied institutional access and legitimacy. \hfill\tinf
  \end{statement}
  \source{Digital Asset launch release, 2016 financing release and current corporate biographies.}
  \note{[Sources] Digital Asset 11 March 2015 release; current about page; 2016 financing announcement. Interpretation is clearly labelled reconstruction (inference tier).}
\end{frame}

\begin{frame}{Technology assembled by acquisition}
  \fontsize{9.9}{11.8}\selectfont
  \begin{tabularx}{\textwidth}{L{1.55cm}L{2.55cm}Y}
    \toprule
    \textbf{Date} & \textbf{Acquisition} & \textbf{Capability added} \\
    \midrule
    June 2015 & Hyperledger Inc. & permissioned ledger, team and trademark \\
    June 2015 & Bits of Proof & Bitcoin and enterprise-ledger engineering \\
    2015/16 & Blockstack.io & private-blockchain services; unrelated to Stacks \\
    April 2016 & Elevence & contractual modelling language that became Daml \\
    2023 onward & Canton network & interoperable private Daml applications \\
    \bottomrule
  \end{tabularx}
  \vspace{0.14cm}
  \muted{\fontsize{9.0}{10.8}\selectfont Hyperledger mark and code were donated to the Linux Foundation (2015--16); Fabric is IBM lineage, Besu is Consensys lineage --- Digital Asset owns neither.}
  \vspace{0.10cm}
  \begin{statement}
    Canton should not be projected backwards into the platform ASX selected in 2016.
  \end{statement}
  \source{Digital Asset corporate milestones; acquisition announcements; Linux Foundation formation and 2016 code-proposal announcements.}
  \note{[Sources] Digital Asset about/milestones; contemporaneous acquisition releases; LF Decentralized Trust 17 Dec 2015 and 9 Feb 2016 announcements. Absorbs the earlier Hyperledger-ownership slide.}
\end{frame}

\begin{frame}{The Ethereum lineage is separate}
  \fontsize{9.8}{11.8}\selectfont
  \begin{tabularx}{\textwidth}{L{2.15cm}Y Y}
    \toprule
    \textbf{Date} & \textbf{Ethereum} & \textbf{Digital Asset} \\
    \midrule
    Nov 2013 & white paper & -- \\
    Jan--Sep 2014 & announced, specified, funded & organised later in 2014 \\
    Mar 2015 & prototypes & public company launch \\
    Jul 2015 & mainnet launch & operating and acquiring teams \\
    \bottomrule
  \end{tabularx}
  \begin{statement}
    Digital Asset predates the Ethereum mainnet, not the Ethereum project. Later overlap came through JPMorgan, Consensys, Hyperledger and shared investors---not a common founding team.
  \end{statement}
  \source{Ethereum official history; Digital Asset March 2015 launch.}
  \note{[Sources] Ethereum.org history; Digital Asset launch release. Merges the earlier chronology and founder-overlap slides; full founder lists remain in the companion report.}
\end{frame}

\begin{frame}{The 2016 coalition changed the company, and kept paying}
  \fontsize{9.6}{11.4}\selectfont
  \begin{columns}[T,onlytextwidth]
    \column{0.50\textwidth}
      \rust{US\$50m+ strategic investors (2016)}
      \begin{itemize}\verytightitems
        \item ABN AMRO, BNP Paribas, Citi, JPMorgan, PNC, Santander
        \item ASX, CME, Deutsche B\"orse, DTCC, ICAP
        \item Accenture and Broadridge
      \end{itemize}
    \column{0.46\textwidth}
      \blue{Board representation}
      \begin{itemize}\verytightitems
        \item BNP Paribas
        \item Deutsche B\"orse
        \item JPMorgan
        \item DTCC
      \end{itemize}
  \end{columns}
  \vspace{0.22cm}
  \begin{tabularx}{\textwidth}{L{1.0cm}L{2.1cm}Y}
    \toprule
    2016 & US\$50m+ & banks, exchanges, CSD/clearing and integrators \\
    2025 & US\$135m & Tradeweb/DRW plus dealers, DTCC and crypto firms \\
    2026 & US\$355m & a16z, ADIA, banks, brokers, exchanges and crypto investors \\
    \bottomrule
  \end{tabularx}
  \vspace{0.14cm}
  \muted{\fontsize{9}{10.8}\selectfont 700+ Canton ecosystem participants reported June 2026; the 2026 round disclosed no valuation.}
  \source{Digital Asset/Broadridge 21 January 2016 financing announcement; 2025 and 11 June 2026 financing releases.}
  \note{[Sources] Digital Asset financing releases. No undisclosed valuation is inferred. Merges the earlier coalition and capital slides. Statement point: customers, infrastructure owners and potential competitors financed and governed the vendor.}
\end{frame}

\begin{frame}{Canton coordinates private applications} % FAST
  \begin{itemize}\tightitems
    \item Participants hold party-specific contract views rather than one public payload ledger.
    \item Synchronisers sequence encrypted messages.
    \item Mediators coordinate atomic commit across applications.
    \item Global Synchronizer governance uses Foundation and super-validator voting.
  \end{itemize}
  \begin{statement}
    Canton distributes data but can still concentrate operational power in synchronisation, validation, hosting and Daml expertise.
  \end{statement}
  \source{Digital Asset Canton participant architecture; Canton Foundation governance.}
  \note{[Sources] Digital Asset Canton documentation; Canton Foundation.}
\end{frame}

\begin{frame}{Institutions hedge every architecture}
  \begin{itemize}\tightitems
    \item JPMorgan: Agor\'a, Swift design, TCH, Kinexys, Partior and Canton.
    \item DTCC: Digital Asset investor; Canton and Besu deployments.
    \item Euroclear: Agor\'a, Fnality and Canton collateral network.
    \item Visa and Mastercard: Agor\'a, Swift/Besu groups and Open USD.
    \item Tradeweb: Digital Asset and Securitize; majority owned by LSEG.
  \end{itemize}
  \begin{statement}
    Strategic optionality lets incumbent institutions shape the winner whichever stack prevails.
  \end{statement}
  \source{Corporate investment, participant and partnership disclosures mapped in the companion report.}
  \note{[Sources] JPMorgan, DTCC, Euroclear, Visa, Mastercard, Tradeweb/LSEG and project disclosures.}
\end{frame}

% ============================ PART VII ============================
\sectionframe{PART VII}{ASX: failure, conflict and consequence}{9 minutes}

\begin{frame}{ASX made the market a proving ground}
  \begin{itemize}\tightitems
    \item January 2016: exploratory Digital Asset work plus 5\% equity stake; June 2016: 8.5\% and a board seat.
    \item 2017: commitment to DLT-based CHESS replacement.
    \item November 2022: project paused; A\$245--255 million written off.
    \item November 2023: TCS BaNCS selected for a staged replacement.
    \item 20 April 2026: Release 1 (clearing) goes live; Release 2 (settlement and subregister) is projected toward the end of the decade.
  \end{itemize}
  \begin{statement}
    Vendor selection, venture investment, product development and critical-system replacement became inseparable.
  \end{statement}
  \source{ASX announcements; ASIC inquiry and Federal Court penalty materials.}
  \note{[Sources] ASX project history; ASIC 2026 inquiry and penalty; ASX Release 1 go-live announcements (April 2026).}
\end{frame}

\begin{frame}{Four risks were accepted together}
  \begin{tabularx}{\textwidth}{L{2.2cm}Y}
    \toprule
    \textbf{Risk} & \textbf{Why it was foreseeable} \\
    \midrule
    Technology & no equivalent national production deployment \\
    Integration & acquired teams and components still forming one stack \\
    Governance & ASX was buyer, investor and board participant \\
    Disclosure & admitting failure damaged both infrastructure and investment narratives \\
    \bottomrule
  \end{tabularx}
  \begin{statement}
    ``Incompetence'' is incomplete: this was a foreseeable institutional-design failure. \hfill\tinf
  \end{statement}
  \source{ASX chronology; Digital Asset acquisition/funding history; ASIC findings.}
  \note{[Sources] ASX, Digital Asset and ASIC materials. Final sentence is an analytical conclusion (inference tier).}
\end{frame}

\begin{frame}{The conflict was real}
  \begin{itemize}\tightitems
    \item ASX had a financial and reputational interest in vendor success.
    \item CME, Deutsche B\"orse, DTCC and major banks invested alongside it.
    \item Those firms gained exposure to reusable post-trade technology and requirements knowledge.
    \item Digital Asset lacked an equivalent live national-market deployment.
  \end{itemize}
  \begin{statement}
    Conflict establishes a duty for exceptional independence. It does not by itself establish corrupt action.
  \end{statement}
  \source{Digital Asset 2016 investor/board release; ASX disclosures; ASIC inquiry.}
  \note{[Sources] Digital Asset/Broadridge financing release; ASX; ASIC.}
\end{frame}

\begin{frame}{The official findings are already severe}
  \begin{itemize}\tightitems
    \item ASIC inquiry: resilience compromised in pursuit of high shareholder returns; governance, capability and stewardship failures identified.
    \item ASX admitted its 10 February 2022 ``progressing well'' statement was misleading --- contravening ss 12DA and 12DB(1)(a) and (e) of the ASIC Act.
    \item July 2026: the Federal Court imposed a A\$20.5 million penalty plus A\$3 million towards ASIC's costs.
  \end{itemize}
  \begin{statement}
    The public record establishes conflicted and misleading stewardship of critical infrastructure. \hfill\tdoc
  \end{statement}
  \source{ASIC, \href{https://www.asic.gov.au/about-asic/news-centre/find-a-media-release/2026-releases/26-059mr-asic-publishes-asx-inquiry-panel-final-report-and-acknowledges-observations/}{April 2026 inquiry findings}; ASIC, \href{https://www.asic.gov.au/about-asic/news-centre/find-a-media-release/2026-releases/26-143mr-asx-ordered-to-pay-20-5-million-penalty-for-misleading-conduct-relating-to-chess-replacement-project/}{July 2026 CHESS penalty}.}
  \note{[Sources] ASIC 26-059MR and 26-143MR; ASIC 26-119MR (admission, statutory provisions, costs order).}
\end{frame}

\begin{frame}{The sabotage hypotheses rank differently}
  \fontsize{9.8}{11.8}\selectfont
  \begin{tabularx}{\textwidth}{L{0.55cm}Y L{3.0cm}}
    \toprule
    & \textbf{Hypothesis} & \textbf{Assessment} \\
    \midrule
    1 & governance failure and shareholder capture & \tdoc \\
    2 & reckless experimentation in conflicted procurement & \tdoc \\
    3 & ecosystem enclosure without central conspiracy & \tinf, partly evidenced \\
    4 & tacit tolerance of failure by beneficiaries & \tinf, weak \\
    5 & coordinated sabotage for a foreign takeover & \tunp \\
    \bottomrule
  \end{tabularx}
  \source{Evidence ledger and counterfactual assessment in the companion report.}
  \note{[Sources] ASIC/ASX record plus documented ownership and chronology. Assessments now use the deck's evidence tiers directly.}
\end{frame}

\begin{frame}{Later benefit is not proof of authorship}
  \begin{itemize}\tightitems
    \item Digital Asset gained requirements knowledge, credibility and a flagship narrative.
    \item Its investors gained option value across post-trade architectures.
    \item ASX sold its remaining 5.4\% stake in 2025 for about A\$57 million.
    \item ASX reported only about A\$10 million pre-tax gain over acquisition cost---far below the project write-off.
  \end{itemize}
  \begin{statement}
    The strategic springboard is credible; a disclosed ASX investment-profit motive is not. \hfill\tinf
  \end{statement}
  \source{ASX 2025 sale announcement and project write-off disclosures.}
  \note{[Sources] ASX 2025 sale announcement and historical project cost disclosures.}
\end{frame}

\begin{frame}{Cboe and TMX did not replace CHESS}
  \begin{itemize}\tightitems
    \item ASIC introduced competition rules limiting misuse of ASX monopoly power.
    \item ASIC approved Cboe Australia as a listing venue in October 2025.
    \item TMX agreed in 2026 to buy Cboe Australia and Canada, subject to approval.
    \item ASX---not Clearstream or TMX---selected TCS BaNCS and remains the CHESS operator.
  \end{itemize}
  \begin{statement}
    De-monopolisation is real; the claim that Clearstream ``swooped in and took CHESS'' is \tunp.
  \end{statement}
  \source{ASIC competition/listing releases; TMX/Cboe transaction; ASX TCS selection.}
  \note{[Sources] ASIC 25-019MR and 25-227MR; TMX/Cboe; ASX.}
\end{frame}

% ============================ PART VIII ============================
\sectionframe{PART VIII}{The test for any jurisdiction}{12 minutes}

\begin{frame}{Dominance without issuing AUD}
  \begin{itemize}\tightitems
    \item Select correspondents and executable FX paths.
    \item Determine collateral eligibility and substitution.
    \item Sequence obligations and reserve liquidity.
    \item Encode identity, sanctions and market-access rules.
    \item Send only the final net or gross instruction to RITS.
  \end{itemize}
  \begin{statement}
    RITS could retain legal finality while becoming the last actuator of decisions made elsewhere.
  \end{statement}
  \source{Agor\'a, Swift, CLS, Meridian FX and synchronisation architectures analysed in the companion report.}
  \note{[Sources] BIS Agorá and Meridian FX; Swift ledger; CLS; central-bank synchronisation work.}
\end{frame}

\begin{frame}{A CLS-style utility can capture workflow}
  \begin{columns}[T,onlytextwidth]
    \column{0.47\textwidth}
      \rust{What Australia keeps}
      \begin{itemize}\tightitems
        \item AUD issuance
        \item ESA balances
        \item RITS finality
        \item domestic emergency law
      \end{itemize}
    \column{0.49\textwidth}
      \blue{What it could lose}
      \begin{itemize}\tightitems
        \item price discovery and routing
        \item collateral allocation
        \item transaction data
        \item operational continuity
      \end{itemize}
  \end{columns}
  \begin{statement}
    Final settlement in public money can coexist with dependence everywhere else.
  \end{statement}
  \note{[Sources] Analytical synthesis based on CLS and proposed orchestration models.}
\end{frame}

\begin{frame}{The case for conceding layers}
  \begin{itemize}\tightitems
    \item \rust{Scale economics}: duplicating clearing, identity and orchestration nationally is expensive; global pools price and net better.
    \item \rust{Capability scarcity}: the RBA's own partly-observed PFMI ratings show that public operation is not automatically resilient.
    \item \rust{Standardisation}: common rails reduce fragmentation, error rates and compliance cost for every participant.
  \end{itemize}
  \vspace{0.26cm}
  \begin{statement}
    The efficiency case is real. It argues for shared layers on contestable terms --- not for uninspectable, unsubstitutable ones. The dispute is about terms, not participation.
  \end{statement}
  \source{RBA 2026 RITS self-assessment; CPMI-IOSCO efficiency principles; companion-report counterarguments chapter.}
  \note{[Sources] RBA RITS ratings; PFMI. Steelman slide: pre-empts the efficiency objection before the policy asks.}
\end{frame}

\begin{frame}{Australia is already piloting the future stack}
  \begin{itemize}\tightitems
    \item \rust{Project Acacia}: RBA and DFCRC wholesale tokenisation programme, spanning the A\$350 billion-plus repo market and other asset classes.
    \item CBA--J.P.\ Morgan (Kinexys)--ASX--HQLAx use case, with Deutsche B\"orse Group as HQLAx's trusted third party; completed May 2026.
    \item \rust{Repo transactions settled using a CBA deposit token} alongside a pilot wholesale CBDC; the securities leg ran on HQLAx with tokenised assets held at ASX.
  \end{itemize}
  \vspace{0.24cm}
  \begin{statement}
    The domestic pilot already contains every element of this briefing: commercial-bank settlement money, a US bank's platform and a Deutsche B\"orse-linked collateral layer above local assets.\\\hspace*{\fill}\tdoc
  \end{statement}
  \source{RBA, \href{https://www.rba.gov.au/media-releases/2024/mr-24-25.html}{Project Acacia}; CBA, \href{https://www.commbank.com.au/articles/newsroom/2026/05/rba-wholesale-funding-markets-trial.html}{May 2026 completion}; J.P.\ Morgan, \href{https://www.jpmorgan.com/payments/newsroom/kinexys-project-acacia-tokenization-trial-australia}{Kinexys trial}.}
  \note{[Sources] RBA/DFCRC Project Acacia; CBA and J.P. Morgan May 2026 releases. Key fact: settlement occurred in commercial-bank money (CBA deposit token) as well as pilot wCBDC --- the settlement-asset question is live domestically, now.}
\end{frame}

\begin{frame}{The tests map onto existing law}
  \fontsize{9.2}{10.9}\selectfont
  \begin{tabularx}{\textwidth}{L{3.25cm}Y}
    \toprule
    \textbf{Instrument} & \textbf{What it reaches} \\
    \midrule
    Payment Systems (Regulation) Act, as amended & ministerial designation reaches orchestration layers and foreign operators \\
    \addlinespace[0.16em]
    RBA financial stability / CS-facility standards & binding conditions on licensed facilities, incl.\ overseas licensees such as LCH \\
    \addlinespace[0.16em]
    APRA CPS 230 & material-service-provider obligations: the home for rehearsed substitution and exit \\
    \addlinespace[0.16em]
    Security of Critical Infrastructure Act & obligations and assistance powers for designated financial assets \\
    \bottomrule
  \end{tabularx}
  \vspace{0.10cm}
  \begin{statement}
    The machinery exists. The unfinished work is applying it deliberately to layers 4--8.
  \end{statement}
  \source{Treasury payments-licensing reforms; RBA standards; APRA CPS 230; SOCI Act.}
  \note{[Sources] Amended PSRA and designation power; RBA FSS/CS-facility standards; APRA CPS 230; SOCI Act. Converts the policy asks from aspiration to mechanism.}
\end{frame}

\begin{frame}{The trusted optimiser must pass seven tests}
  \fontsize{10.3}{12.4}\selectfont
  \begin{enumerate}\tightitems
    \item Can the RBA inspect and reproduce every material decision?
    \item Are liquidity locks participant-authorised and bounded?
    \item Can the RBA revoke them during an emergency?
    \item Is there a direct route that bypasses the optimiser?
    \item Can another operator substitute without data loss?
    \item Does Australia retain keys, telemetry and recovery authority?
    \item Can the AUD domain isolate and continue domestically?
  \end{enumerate}
  \note{[Sources] Control tests developed in the companion report.}
\end{frame}

\begin{frame}{The settlement rule must be explicit}
  \begin{statement}
    A net interbank obligation denominated in Australian dollars is finally discharged only in RBA money, unless the RBA expressly authorises another settlement asset for a defined use case.
  \end{statement}
  \vspace{0.42cm}
  \begin{itemize}\tightitems
    \item Token movement is not enough.
    \item Contractual netting is not enough.
    \item A foreign operator's commitment record is not enough.
  \end{itemize}
  \note{[Sources] Policy recommendation from The Sovereign Ledger. Lands directly after Acacia: the pilot's commercial-bank-money settlement is exactly the case the rule must anticipate.}
\end{frame}

\begin{frame}{External ESA locks need emergency limits}
  \begin{statement}
    External ESA locks must be capped, participant-authorised and revocable under RBA emergency rules.
  \end{statement}
  \vspace{0.42cm}
  \begin{itemize}\tightitems
    \item Maximum amount and duration set ex ante
    \item Explicit participant consent for every reservation class
    \item RBA visibility and unilateral crisis release authority
    \item No permanent exclusive synchroniser
    \item Tested fallback to direct RITS submission
  \end{itemize}
  \note{[Sources] Policy recommendation from The Sovereign Ledger.}
\end{frame}

\begin{frame}{Continuity requires rehearsed exit}
  \begin{itemize}\tightitems
    \item Domestic copies of authoritative data and audit trails
    \item Source, build, key and privileged-access arrangements
    \item Operator substitution and manual contingency modes
    \item Publicly funded engineering capability inside Australia
    \item Crisis exercises that assume the external layer is unavailable
  \end{itemize}
  \begin{statement}
    Exit that has never been rehearsed is not an exit plan.
  \end{statement}
  \note{[Sources] Policy synthesis from the companion report and PFMI operational-resilience principles.}
\end{frame}

\begin{frame}{The conclusion is adversarial, not conspiratorial}
  \claim{Retaining the last ledger entry is insufficient if private institutions control every condition that precedes it.}
  \vspace{0.30cm}
  \begin{itemize}\tightitems
    \item Assume every actor will exploit lawful control and switching costs.
    \item Demand transparent optimisation and public emergency authority.
    \item Preserve multiple routes, domestic competence and credible substitution.
    \item Distinguish documented conflict from unproved coordination.
  \end{itemize}
  \source{The Sovereign Ledger integrated briefing, 5 August 2026.}
  \note{[Sources] Synthesis of the full companion report.}
\end{frame}

% ============================ CLOSE ============================
\begingroup
\setbeamercolor{background canvas}{bg=Ink}
\begin{frame}[plain]
  \vspace{0.95cm}
  {\color{Gold}\fontsize{10}{12}\selectfont\bfseries FINAL PRINCIPLE\par}
  \vspace{0.45cm}
  {\color{White}\fontsize{24}{28}\selectfont\bfseries
    No private owner, foreign parent, bank club or technology supplier should be able to hold the continuity of money and title hostage.\par}
  \vfill
  {\color{Mist}\fontsize{10.3}{12.8}\selectfont Contestable control is the operational meaning of sovereignty.}
  \note{[Sources] Closing conclusion of The Sovereign Ledger. The single deliberate use of the word.}
\end{frame}
\endgroup

% ============================ APPENDIX ============================
\appendix
\sectionframe[Reference material]{APPENDIX}{Participants, ownership and sources}{not presented}

\begin{frame}{Agor\'a participants: global banks}
  \fontsize{10}{12}\selectfont
  \begin{columns}[T,onlytextwidth]
    \column{0.48\textwidth}
      \begin{itemize}\verytightitems
        \item BNP Paribas
        \item BNY
        \item Citi
        \item Cr\'edit Agricole CIB
        \item Deutsche Bank
        \item HSBC
        \item JPMorgan Chase
        \item Standard Chartered
        \item UBS
      \end{itemize}
    \column{0.48\textwidth}
      \begin{itemize}\verytightitems
        \item Santander
        \item BBVA
        \item CaixaBank
        \item Commerzbank
        \item Groupe BPCE
        \item Lloyds
        \item NatWest
        \item FNBO
        \item TD Bank N.A.
      \end{itemize}
  \end{columns}
  \source{BIS/IIF Agor\'a final participant list.}
  \note{[Sources] BIS/IIF participant list.}
\end{frame}

\begin{frame}{Agor\'a participants: Japan and Korea}
  \fontsize{10.3}{12.4}\selectfont
  \begin{columns}[T,onlytextwidth]
    \column{0.46\textwidth}
      \rust{Japan}
      \begin{itemize}\verytightitems
        \item Mizuho
        \item MUFG
        \item SBI Shinsei
        \item SMBC
      \end{itemize}
    \column{0.50\textwidth}
      \blue{Korea}
      \begin{itemize}\verytightitems
        \item Hana
        \item Industrial Bank of Korea
        \item KB Kookmin
        \item NongHyup
        \item Shinhan
        \item Woori
      \end{itemize}
  \end{columns}
  \source{BIS/IIF Agor\'a final participant list.}
  \note{[Sources] BIS/IIF participant list.}
\end{frame}

\begin{frame}{Agor\'a participants: specialist banks}
  \fontsize{10}{12}\selectfont
  \begin{columns}[T,onlytextwidth]
    \column{0.48\textwidth}
      \rust{Mexico / Latin America}
      \begin{itemize}\verytightitems
        \item Banorte
        \item Monex
        \item Banco BV
      \end{itemize}
      \vspace{0.22cm}
      \rust{Switzerland}
      \begin{itemize}\verytightitems
        \item AMINA
        \item BCV
        \item Basler Kantonalbank
      \end{itemize}
    \column{0.48\textwidth}
      \blue{Switzerland continued}
      \begin{itemize}\verytightitems
        \item PostFinance
        \item Sygnum
      \end{itemize}
      \vspace{0.22cm}
      \muted{These institutions add regional, crypto-bank, cooperative and specialist-FX perspectives.}
  \end{columns}
  \source{BIS/IIF Agor\'a final participant list.}
  \note{[Sources] BIS/IIF participant list.}
\end{frame}

\begin{frame}{Agor\'a participants: FMIs and networks}
  \begin{itemize}\tightitems
    \item \rust{Eurex Clearing} --- sole CCP
    \item \rust{Euroclear} --- CSD/ICSD, custody and collateral
    \item \rust{SIX Digital Exchange} --- regulated DLT CSD/exchange
    \item \rust{Mastercard} --- network rules, identity and distribution
    \item \rust{Visa} --- network/risk rules and bank-token platform
    \item \rust{Swift} --- ISO 20022, connectivity and competing ledger strategy
  \end{itemize}
  \source{BIS/IIF participant list; participant corporate disclosures.}
  \note{[Sources] BIS/IIF list; Eurex, Euroclear, SIX, Mastercard, Visa and Swift.}
\end{frame}

\begin{frame}{Other projects test different control points}
  \fontsize{9.8}{11.8}\selectfont
  \begin{tabularx}{\textwidth}{L{1.95cm}Y Y}
    \toprule
    \textbf{Project} & \textbf{Mechanism} & \textbf{Control question} \\
    \midrule
    Mariana & AMM for wCBDC FX & who sets curves and supplies liquidity? \\
    Cedar--Ubin+ & atomic links across separate ledgers & can interoperability avoid a central platform? \\
    Helvetia & live wCBDC versus RTGS link & integrated token or synchronised conventional money? \\
    Pontes & DLT-to-TARGET bridge & how does market DLT reach public money? \\
    Meridian FX & external queue and liquidity saving & can an optimiser shape settlement before RTGS? \\
    \bottomrule
  \end{tabularx}
  \source{BIS Innovation Hub project reports; ECB Pontes; SNB Helvetia; New York Fed Cedar.}
  \note{[Sources] BIS Mariana and Meridian; NY Fed Cedar; ECB Pontes; SNB Helvetia. Moved from Part IV.}
\end{frame}

\begin{frame}{Ownership atlas: incumbent infrastructure}
  \fontsize{8.8}{10.5}\selectfont
  \begin{tabularx}{\textwidth}{L{1.65cm}Y Y}
    \toprule
    \textbf{Institution} & \textbf{Ownership} & \textbf{Control point} \\
    \midrule
    Swift & financial-institution cooperative & messaging and future ledger \\
    CLS & large financial institutions & FX acceptance and netting \\
    TCH & large US member banks & CHIPS, RTP and token proposal \\
    DTCC & user/industry owned & US post-trade chain \\
    Clearstream & Deutsche B\"orse subsidiary & CSD/ICSD and collateral \\
    Euroclear & mixed private, public and sovereign holders & CSD/ICSD and collateral \\
    ASX & listed company & Australian clearing and CHESS \\
    \bottomrule
  \end{tabularx}
  \note{[Sources] Corporate governance and shareholder disclosures in companion report.}
\end{frame}

\begin{frame}{Ownership atlas: emerging infrastructure}
  \fontsize{8.8}{10.5}\selectfont
  \begin{tabularx}{\textwidth}{L{1.75cm}Y Y}
    \toprule
    \textbf{Institution} & \textbf{Ownership / governance} & \textbf{Control point} \\
    \midrule
    Digital Asset & private strategic/venture coalition & Daml/Canton software and ecosystem \\
    Canton Foundation & member/super-validator voting & Global Synchronizer governance \\
    Fnality & 24 banks/FMIs & reserve-backed Besu payment systems \\
    R3 & private strategic investors & Corda product and expertise \\
    Consensys & private founder/venture group & Ethereum tools, Infura, Besu lineage \\
    Fireblocks & private venture-backed vendor & keys, policy engine and nodes \\
    Securitize & NYSE-listed SECZ & KYC, transfer agent, ATS and funds \\
    \bottomrule
  \end{tabularx}
  \note{[Sources] Corporate and foundation disclosures cited in companion report.}
\end{frame}

\begin{frame}{Primary sources: public-money architecture}
  \fontsize{8.1}{9.7}\selectfont
  \begin{itemize}\verytightitems
    \item \href{https://www.rba.gov.au/payments-and-infrastructure/rits/about.html}{Reserve Bank of Australia --- RITS}
    \item \href{https://www.bankofengland.co.uk/payment-and-settlement/a-brief-introduction-to-the-real-time-gross-settlement-system-and-chaps}{Bank of England --- RTGS and CHAPS}
    \item \href{https://www.ecb.europa.eu/paym/target/html/index.en.html}{ECB --- TARGET Services}
    \item \href{https://www.snb.ch/en/the-snb/mandates-goals/payment-transactions/projekt_helvetia}{SNB --- Project Helvetia}
    \item \href{https://www.payments.ca/systems-services/payment-systems/high-value-payment-system-lynx}{Payments Canada --- Lynx}
    \item \href{https://www.bis.org/publ/othp110.pdf}{BIS --- Project Agor\'a final report}
  \end{itemize}
  \note{[Sources] Links shown on slide.}
\end{frame}

\begin{frame}{Primary sources: projects and ledgers}
  \fontsize{8.1}{9.7}\selectfont
  \begin{itemize}\verytightitems
    \item \href{https://www.bis.org/publ/othp47.htm}{BIS --- Project Dunbar}
    \item \href{https://www.bis.org/publ/othp44.htm}{BIS --- Project Jura}
    \item \href{https://www.bis.org/about/bisih/topics/cbdc/mcbdc_bridge.htm}{BIS --- Project mBridge}
    \item \href{https://www.bis.org/about/bisih/topics/fmis/agora/rvt.htm}{BIS --- Agor\'a real-value testing}
    \item \href{https://www.newyorkfed.org/aboutthefed/nyic/project-cedar}{New York Fed --- Project Cedar}
    \item \href{https://www.swift.com/news-events/press-releases/swifts-blockchain-ledger-ready-use-17-banks-set-pioneer-tokenised-cross-border-payments-trusted-global-infrastructure}{Swift --- 17-bank ledger cohort}
    \item \href{https://www.theclearinghouse.org/payment-systems/Articles/2026/06/Major-Financial-Institutions-Unveil-Bank-Led-On-Chain-Money-Initiative}{TCH --- bank-led on-chain money}
  \end{itemize}
  \note{[Sources] Links shown on slide.}
\end{frame}

\begin{frame}{Primary sources: technology lineage}
  \fontsize{8.1}{9.7}\selectfont
  \begin{itemize}\verytightitems
    \item \href{https://www.prnewswire.com/news-releases/digital-asset-holdings-delivers-first-of-its-kind-technology-to-bridge-the-gap-between-crypto-and-mainstream-assets-300048767.html}{Digital Asset --- March 2015 launch}
    \item \href{https://www.digitalasset.com/about}{Digital Asset --- milestones and founders}
    \item \href{https://www.broadridge.com/press-release/2016/digital-asset-closes-funding-round-exceeding-50-million}{Digital Asset --- 2016 institutional round}
    \item \href{https://www.lfdecentralizedtrust.org/announcements/2015/12/17/linux-foundation-unites-industry-leaders-to-advance-blockchain-technology}{Linux Foundation --- Hyperledger formation}
    \item \href{https://www.lfdecentralizedtrust.org/projects/besu}{LF Decentralized Trust --- Besu}
    \item \href{https://ethereum.org/en/ethereum-history-founder-and-ownership/}{Ethereum --- official history}
  \end{itemize}
  \note{[Sources] Links shown on slide.}
\end{frame}

\begin{frame}{Primary sources: ASX and ownership}
  \fontsize{8.1}{9.7}\selectfont
  \begin{itemize}\verytightitems
    \item \href{https://www.asic.gov.au/about-asic/news-centre/find-a-media-release/2026-releases/26-059mr-asic-publishes-asx-inquiry-panel-final-report-and-acknowledges-observations/}{ASIC --- ASX inquiry findings}
    \item \href{https://www.asic.gov.au/about-asic/news-centre/find-a-media-release/2026-releases/26-143mr-asx-ordered-to-pay-20-5-million-penalty-for-misleading-conduct-relating-to-chess-replacement-project/}{ASIC --- CHESS penalty}
    \item \href{https://www.swift.com/about-us/organisation-governance/swift-governance}{Swift --- cooperative governance}
    \item \href{https://www.cls-group.com/about/shareholders/}{CLS --- shareholders}
    \item \href{https://www.euroclear.com/investorrelations/en/top-10-shareholders-Euroclear-Holding-SA-NV.html}{Euroclear --- major shareholders}
    \item \href{https://deutsche-boerse.com/dbg-en/media/news-stories/press-releases/Clearstream-Unveils-Next-Generation-Digital-Securities-Infrastructure-5318978}{Deutsche B\"orse / Clearstream}
  \end{itemize}
  \note{[Sources] Links shown on slide.}
\end{frame}

\begin{frame}{Primary sources: Australia and clearing}
  \fontsize{8.1}{9.7}\selectfont
  \begin{itemize}\verytightitems
    \item \href{https://www.rba.gov.au/media-releases/2024/mr-24-25.html}{RBA --- Project Acacia}
    \item \href{https://www.commbank.com.au/articles/newsroom/2026/05/rba-wholesale-funding-markets-trial.html}{CBA --- Acacia repo trial completion (May 2026)}
    \item \href{https://www.jpmorgan.com/payments/newsroom/kinexys-project-acacia-tokenization-trial-australia}{J.P.\ Morgan --- Kinexys Acacia trial}
    \item \href{https://www.cfr.gov.au/publications/policy-statements-and-other-reports/2013/report-on-the-australian-otc-derivatives-market-july/}{CFR --- Australian OTC derivatives market reports}
    \item \href{https://www.lseg.com/en/post-trade/clearing/lch-services/swapclear}{LSEG --- LCH SwapClear}
    \item \href{https://www.indeval.com.mx/en/product?p=clearing-settlement}{Indeval --- DAL\'I clearing and settlement}
  \end{itemize}
  \note{[Sources] Links shown on slide.}
\end{frame}

\end{document}
